\documentclass[11pt]{article}
\usepackage[margin=1in]{geometry}
\usepackage[utf8]{inputenc}
\usepackage[T1]{fontenc}
\usepackage{lmodern}
\usepackage{amsmath,amssymb}
\usepackage{graphicx}
\usepackage{booktabs}
\usepackage{array}
\newcolumntype{L}[1]{>{\raggedright\arraybackslash}p{#1}}
\usepackage{xcolor}
\usepackage{listings}
\usepackage[hidelinks]{hyperref}
\usepackage{url}
\usepackage{caption}
\captionsetup{font=small,labelfont=bf}

\newcommand{\memprint}{MemPrint}
\newcommand{\secref}[1]{Section~\ref{#1}}
\newcommand{\figref}[1]{Figure~\ref{#1}}
\newcommand{\tabref}[1]{Table~\ref{#1}}
\newcommand{\term}[1]{\emph{#1}}
\makeatletter
\newcommand{\tabinput}[1]{\@@input #1 }   % \input that a table row may start with \multicolumn
\makeatother
\lstset{basicstyle=\ttfamily\footnotesize,columns=fullflexible,keepspaces=true,frame=none,
        xleftmargin=1em,showstringspaces=false}

\title{Memory Footprints from Resident Memory and Allocator Statistics:\\
A Footprint Model at the Cost of a Native Run}
\author{Nafis Mustakin\\University of California, Riverside}
\date{October 2026}

\begin{document}
\setlength{\emergencystretch}{3em}
\maketitle

\begin{abstract}
A memory-aware scheduler needs to know how much memory a job uses over its run. \memprint{}, a Pin-based
footprint profiler, measures the pages a program touches, but its windowed mode costs $12.0\times$ a native run at
the median of 22 programs, $18.7\times$ to $93.0\times$ on programs that allocate and free memory often, and up to
$67.4\times$ on a language model engine. The resident set size (RSS) of a native run costs nothing extra, but it
overestimates allocator-heavy programs by a median of 16\% over the whole run, and miniVite by 102\%, because the
allocator keeps freed memory resident and RSS counts shared libraries. We propose a footprint model that samples, from
one thread inside the process every 5\,ms, the anonymous RSS, the proportional set size, the RSS of mapped data files,
and the free bytes the allocator reports (glibc, jemalloc or TCMalloc). Its data estimate, anonymous RSS minus free
bytes plus mapped data files and the runtime's pages, has a median whole-run error of $-4\%$ on the 7 allocator-heavy
programs of a 23-program corpus ($-7\%$ on miniVite), against $-5\%$ for \memprint{}, at a median of $0.99\times$ the
run time of a plain RSS run. Transferred from smaller inputs to a held-out larger one, it has a median error of
$-4\%$, against $+20\%$ for RSS and $-6\%$ for \memprint{}. In a cluster scheduling simulator that reserves memory by
each job's estimated curve, it gives the lowest mean response time of all estimates that profile the job's own input
(433 minutes at a load of 0.8, against 515 for the RSS curve and 519 for \memprint{}). Under TCMalloc its median
error is $-4\%$. Under jemalloc it is $-1\%$ with transparent huge pages disabled and $+35\%$ with them enabled,
because the kernel backs memory the program touches in part with whole 2\,MiB pages. On llama2.c and llama.cpp, the
data-file term brings the model from between $-39\%$ and $-92\%$ to between $-28\%$ and $+15\%$ of the truth.
\end{abstract}

\section{Introduction}
\label{sec:intro}

A scheduler that packs jobs onto nodes by their memory use needs an estimate of each job's footprint over time,
because a job's peak lasts only part of its run~\cite{KSplus2024,Zacarias2021Accurate,Li2023Pond}. A reservation
that follows the footprint over time lets a node hold more jobs than a reservation of the peak, but the estimate
must not fall below what the job uses, or the job runs out of memory.

\memprint{} measures a program's footprint with Pin~\cite{Luk2005Pin}. It records the bytes and pages each program
touches and the time at which each block is freed. Its windowed mode watches 5\% of the run and costs $12.0\times$ a
native run at the median on 22 programs of our corpus. However, on the programs where a footprint estimate helps the
most, those that allocate and free memory throughout their run, the cost rises to $18.7\times$ to $93.0\times$,
because Pin instruments every memory access and every allocator call.

The resident set size of a native run is the obvious cheap alternative. On array programs, RSS is the footprint,
because these programs allocate their arrays once and touch them. On allocator-heavy programs, RSS overestimates the
footprint over the run by a median of 16\%, and by 102\% on miniVite~\cite{minivite}, because the kernel counts
memory that the program no longer uses.

In this work, we make the following contributions:
\begin{itemize}
\item We identify the parts of RSS that are not part of the footprint, freed memory that the allocator keeps
resident and the pages of shared libraries, and we measure their size per program (\secref{sec:gap}).
\item We propose a footprint model that combines RSS, PSS, the RSS of mapped data files and the allocator's free
bytes, sampled every 5\,ms from a thread inside the process through an \texttt{LD\_PRELOAD} library
(\secref{sec:model}).
\item We show on 23 programs that the model matches \memprint{}'s accuracy on the job's own input and on larger
inputs predicted from smaller ones, at a median of $0.99\times$ the cost of a plain RSS run
(\secref{sec:repeat} and \secref{sec:transfer}).
\item We show in a cluster scheduling simulator that the model's curve gives the lowest response time of all
estimates of the job's own input (\secref{sec:sched}).
\item We evaluate the model under jemalloc and TCMalloc and on two language model engines, and we show that
transparent huge pages are the remaining source of disagreement with a 4\,KiB pages-touched truth
(\secref{sec:alloc} and \secref{sec:llm}).
\item We test four other ways to correct RSS, each alone, and show why each is less accurate
(\secref{sec:options}).
\item We use the model at run time in libmpr, our allocation-resizing runtime, to return glibc's free memory to the
kernel from \texttt{free}, which lowers the time-averaged RSS of the allocator-heavy programs by 4.4\% at a median
run time of $1.01\times$, half of what returning every free byte would give (\secref{sec:trim}).
\item We compare the costs and effects of the approaches to the gap between RSS and the footprint, from estimating it
to lowering it and reserving less than the request (\secref{sec:space}).
\end{itemize}

\section{Background: what RSS counts}
\label{sec:background}

The RSS of a process is the number of its pages in memory. An anonymous page enters RSS when the program first
touches it and leaves when the program or its allocator unmaps it or releases it with \texttt{madvise}. A file-backed
page, such as a page of a shared library's code or of a mapped data file, enters RSS when any process reads it and
stays as long as it is mapped and the kernel does not reclaim it. The \term{footprint} we estimate is the set of
pages the program has touched and not yet freed, plus the pages the runtime needs before the program starts. RSS
differs from the footprint in five ways.

\textbf{Allocator retention.} The C library does not return every freed block to the kernel. glibc's
\texttt{malloc} serves blocks below a threshold (128\,KiB, raised dynamically up to 32\,MiB) from its arenas, keeps
freed blocks in free lists for reuse, and releases memory only from the top of the heap when the free space there
exceeds the trim threshold~\cite{glibcmanual}. jemalloc and TCMalloc also release freed memory late or in
batches~\cite{Evans2006jemalloc,Hunter2021Temeraire}. A program that frees and allocates memory throughout its run
therefore holds freed pages in RSS. Each of these allocators reports the bytes it holds free through its own
interface.

\textbf{Shared libraries.} RSS counts every resident page of every mapped file, including the code of the C library,
the C++ runtime and, for MPI programs, the MPI library. These pages do not depend on the input, and other processes on
the node map the same pages. The PSS of a process divides each resident page by the number of processes that map it.

\textbf{Mapped data files.} A program that maps its input, such as a language model engine that maps its weights,
holds that input in file-backed pages. These pages are part of the footprint, unlike the pages of libraries.

\textbf{Language runtimes with their own allocators.} CPython's object allocator takes 256\,KiB arenas from the system
and returns an arena only when all its objects are freed. GCC's garbage-collected heap also maps its own pages. The C
library's statistics do not cover these allocators.

\textbf{Page size.} RSS counts whole pages. On our corpus, the 4\,KiB pages a program touches exceed the bytes it
touches by 1\% at the median, so we use the pages as the truth. With transparent huge pages~\cite{LinuxTHP}, the
kernel can back a 2\,MiB aligned range of anonymous memory with one huge page at the first touch, and RSS then counts
2\,MiB even if the program touches only part of the range~\cite{Kwon2016Ingens,Panwar2019HawkEye}. Our machine
enables transparent huge pages for all anonymous memory (the \texttt{always} setting).

\section{The footprint model}
\label{sec:model}

The model estimates a footprint curve from five quantities sampled at the same time inside the process. Let $A(t)$
be the anonymous RSS, $P(t)$ the PSS, $M(t)$ the RSS of mapped data files, $F(t)$ the bytes the allocator holds
free, and $B$ the runtime's own pages, the smallest RSS of an earlier run of the program, which is measured once per
program. The model gives two estimates.

\textbf{Estimate 1, the data footprint,} is
\begin{equation}
\hat{f}_1(t) = \max\bigl(A(t) - F(t),\,0\bigr) + M(t) + B .
\end{equation}
It counts the anonymous pages the allocator has not freed and the pages of mapped data files, and it leaves out the
pages of libraries. It estimates the memory that grows with the input, which is what a scheduler predicts and what
our truth measures.

\textbf{Estimate 2, the charged footprint,} is
\begin{equation}
\hat{f}_2(t) = \max\bigl(P(t) - F(t),\,0\bigr) .
\end{equation}
PSS is the anonymous RSS plus each file-backed page divided by the number of processes that map it, so the estimate
charges the job for its share of shared libraries and data files and for all of a library that only it maps. It
estimates the memory the job adds to a node.

Both estimates subtract free bytes from resident memory rather than count live allocations. A counter of live
\texttt{malloc} blocks misses every allocation that does not go through \texttt{malloc}, so an estimate built from it
falls below the truth whenever a program uses its own allocator. Subtracting free bytes errs the other way. Memory
that the allocator does not report stays in RSS, so a program whose allocator the model cannot see gets the RSS
estimate, which is too high but safe.

\begin{sloppypar}
\textbf{Free bytes.} The library looks up the allocator's interface when it starts. With jemalloc (\texttt{mallctl}),
the free bytes are \texttt{stats.resident} minus \texttt{stats.allocated} minus \texttt{stats.metadata}, after a
refresh of the statistics through \texttt{epoch}. With TCMalloc (\texttt{MallocExtension\_GetNumericProperty}), they
are \texttt{generic.heap\_size} minus \texttt{tcmalloc.pageheap\_unmapped\_bytes} minus
\texttt{generic.current\_allocated\_bytes}. Otherwise they are glibc's \texttt{mallinfo} field \texttt{fordblks}, the
free bytes of all its arenas. The machine runs glibc 2.28, which has no \texttt{mallinfo2}, so glibc's free bytes are
a 32-bit field that wraps above 4\,GiB. A newer glibc reports them in 64 bits.
\end{sloppypar}

\textbf{Mapped data files.} A data file is a file mapping that is not the executable, a shared object (a path that
contains \texttt{.so}), or the C library's locale and character set data. $M(t)$ is the sum of the RSS of these
mappings.

\textbf{Implementation.} A shared library of 162 lines of C, loaded with \texttt{LD\_PRELOAD}, starts one thread in
its constructor. Every 5\,ms the thread reads the process's memory counters, queries the allocator, and writes one line
to a buffered file. It has two modes. The first reads \texttt{/proc/self/smaps\_rollup}, which gives $A$, $P$ and RSS
but not $M$. The second reads \texttt{/proc/self/smaps}, one record per mapping, which also gives $M$. On the 23
corpus programs, $M$ is at most 0.07\,MiB, so we evaluate the corpus with the first mode and the language model
engines with the second. The thread does not stop the program's threads, and it takes the allocator's locks only for
the duration of its query. The library removes itself from \texttt{LD\_PRELOAD} so that child processes are not
sampled, and a program made of several processes would need one output file per process.

\section{Methodology}
\label{sec:method}

\textbf{Machine.} All runs use a two-socket Intel Xeon E5-2620 v4 (16 cores, 32 threads) with 62\,GiB of memory,
Linux 4.18 with transparent huge pages enabled, and glibc 2.28. Runs whose time we report execute one at a time on an
otherwise idle machine.

\textbf{Corpus.} We use the 23 programs of the corpus for which a pages-touched truth exists, each at its largest
input: PolyBench (2mm, atax, gemm, jacobi-2d), the NAS parallel benchmarks (cg, is, mg), the GAP benchmark suite (bfs,
cc, pr, sssp)~\cite{beamer2015gap}, the proxy applications HPCCG, Kripke, LULESH, miniFE, miniVite and Quicksilver,
darknet AlexNet inference, GCC's \texttt{cc1} compiling SQLite, and Lua~\cite{ierusalimschy1996lua}, CPython, Perl and
SQLite running table, JSON, hash and SQL workloads. We call 7 programs \term{allocator-heavy}, because they allocate
and free memory throughout their run: miniVite, CPython, Lua, SQLite, Perl, Quicksilver and \texttt{cc1}. We chose the
set before these experiments, from the earlier observation that RSS and \memprint{} disagree on these programs and
agree on the others. The runs last 0.1 to 12\,s and peak at 27 to 273\,MiB of RSS. Every program except darknet also
has one or two smaller inputs. Separately, we use llama2.c with three model sizes (15M, 42M and 110M parameters) and
llama.cpp with a 135M-parameter model generating 32 and 128 tokens.

\textbf{Truth.} The truth is the number of 4\,KiB pages the program touches and has not freed, which Pin measures in
\memprint{}'s splitter mode with page counting and free tracking, plus the runtime's pages $B$. Pin's time axis is
its count of memory references. We place the truth on the wall clock of a native run through hardware counts of the
loads and stores that the native run executes in every 5\,ms interval. The truth at a share $f$ of Pin's references is
placed at the time at which the native run had executed a share $f$ of its references. We divide each run into 100
equal steps of progress and take the largest value of each curve in each step. The estimates are placed on the same
100 steps by the wall clock of their own run. The truths under jemalloc and TCMalloc come from Pin runs with the same
allocator, and they are placed on the wall clock by the counts of the native run with glibc.

\textbf{Metrics.} We report three errors of an estimate $e$ against the truth $t$ over the 100 steps. The signed
whole-run error is $(\sum e - \sum t)/\sum t$, which is positive when the estimate reserves more memory over the run
than the program uses. The absolute whole-run error is $\sum |e - t| / \sum t$. The peak error is $\max e / \max t -
1$. Each program runs 3 times, and we report the median of each error over the runs. The overhead of an estimate is
the median wall time of its run over the median wall time of a run under the plain RSS sampler of the corpus
(\texttt{rss\_sampler.py}, an external process that reads \texttt{smaps\_rollup} every 5\,ms and takes 40\,ms to
start).

\textbf{Baselines.} \emph{RSS} is the RSS curve of the plain sampler's run. \emph{\memprint{}} is the windowed mode
that watches 5\% of the run (W5), placed on the native wall clock by the same reference counts as the truth, plus $B$.
Because the truth and \memprint{} pass through the same clock mapping, \memprint{}'s timing error against the truth
is smaller than its timing error against the real run. Its cost is the wall time of the W5 run and of the counter run
over the wall time of an earlier plain RSS run. \memprint{} has no W5 curve for darknet, so its rows cover 22
programs.

\section{Accuracy on the job's own input}
\label{sec:repeat}

\begin{table}[t]
\centering
\small
\caption{Medians over programs, each the median of 3 runs. Whole-run errors and peak errors are percentages of the
truth (pages touched plus the runtime's pages, on the native clock). The 90th percentile of the absolute error and
the lowest peak error show the worst cases. Overhead is the run time over a plain RSS run ($\times$).}
\label{tab:summary}
\begin{tabular}{llrrrrrrr}
\toprule
 & & & \multicolumn{3}{c}{Whole-run error (\%)} & \multicolumn{2}{c}{Peak error (\%)} & Overhead \\
\cmidrule(lr){4-6}\cmidrule(lr){7-8}
Programs & Estimate & $n$ & signed & absolute & q90 abs. & median & lowest & ($\times$) \\
\midrule
Allocator-heavy (7) & RSS & 7 & +16 & 17 & 55 & +10 & 0 & 1.00\\
 & Data (1) & 7 & $-$4 & 8 & 10 & $-$2 & $-$4 & 0.99\\
 & Shared (2) & 7 & 0 & 9 & 23 & $-$1 & $-$3 & 0.99\\
 & \memprint{} & 7 & $-$5 & 5 & 17 & $-$5 & $-$16 & 34.4\\
\addlinespace
Others (16) & RSS & 16 & +1 & 2 & 8 & 0 & $-$4 & 1.00\\
 & Data (1) & 16 & $-$1 & 3 & 8 & $-$1 & $-$5 & 0.98\\
 & Shared (2) & 16 & $-$5 & 7 & 10 & $-$4 & $-$9 & 0.98\\
 & \memprint{} & 15 & $-$1 & 1 & 3 & $-$1 & $-$2 & 8.9\\
\addlinespace
All (23) & RSS & 23 & +4 & 7 & 20 & +1 & $-$4 & 1.00\\
 & Data (1) & 23 & $-$1 & 5 & 9 & $-$1 & $-$5 & 0.99\\
 & Shared (2) & 23 & $-$4 & 8 & 12 & $-$3 & $-$9 & 0.99\\
 & \memprint{} & 22 & $-$1 & 2 & 12 & $-$1 & $-$16 & 12.0\\

\bottomrule
\end{tabular}
\end{table}

\begin{table}[t]
\centering
\small
\caption{Whole-run error per program (\% of the truth), signed and absolute, median of 3 runs. Overhead is the run
time over a plain RSS run ($\times$). The model's two estimates come from the same run. \memprint{} has no W5 run of
darknet.}
\label{tab:model}
\setlength{\tabcolsep}{4.5pt}
\begin{tabular}{lrrrrrrrrrr}
\toprule
 & \multicolumn{2}{c}{RSS} & \multicolumn{2}{c}{Data (1)} & \multicolumn{2}{c}{Shared (2)} &
 \multicolumn{2}{c}{\memprint{}} & \multicolumn{2}{c}{Overhead ($\times$)} \\
\cmidrule(lr){2-3}\cmidrule(lr){4-5}\cmidrule(lr){6-7}\cmidrule(lr){8-9}\cmidrule(lr){10-11}
Program & sign. & abs. & sign. & abs. & sign. & abs. & sign. & abs. & model & \memprint{} \\
\midrule
\tabinput{tables/model.tex}
\bottomrule
\end{tabular}
\end{table}

\textbf{Allocator-heavy programs.} The data estimate reduces the median whole-run error of RSS from $+16\%$ to $-4\%$
and the median absolute error from 17\% to 8\% (\tabref{tab:summary}). \memprint{} has a median signed error of
$-5\%$ and a median absolute error of 5\%. Per program (\tabref{tab:model}), the data estimate reduces miniVite's
error from $+102\%$ to $-7\%$, CPython's from $+25\%$ to $-3\%$, and Lua's from $+20\%$ to $-1\%$, and its absolute
error is at most 11\% on every allocator-heavy program. Its peak is within 4\% of the truth's on all 7 programs, with a
median of $-2\%$, while the RSS peak is 10\% too high at the median. \memprint{}'s lowest peak error is $-16\%$, on
Lua.

\textbf{The charged estimate.} The charged estimate has a median signed error of 0\% on the allocator-heavy programs,
and it lies above the truth on miniVite ($+37\%$), \texttt{cc1} ($+12\%$) and Quicksilver ($+8\%$), because it keeps
the pages of libraries that only this job maps. Unlike the data estimate, the charged estimate counts memory that the
truth does not count, so its error against the truth does not measure its accuracy as an estimate of a node's memory
use.

\textbf{Other programs.} On the 16 other programs, RSS is already accurate, with a median signed error of $+1\%$ and a
median absolute error of 2\%. The data estimate has $-1\%$ and 3\%, and \memprint{} has $-1\%$ and 1\%.

\textbf{The truth's clock.} An earlier version of this comparison kept the truth on Pin's count of memory references.
Every estimate then had whole-run errors of 10\% or more in the same direction on 7 of the 16 other programs, from
$+10\%$ to $+35\%$ on the GAP programs and from $-10\%$ to $-25\%$ on atax, NPB-IS and NPB-MG, with peak errors within
6\%. A phase that executes few memory references per second, such as GAP's graph construction, takes a smaller share
of the run on Pin's clock than on the wall clock. Placing the truth on the native wall clock removes these errors, and
the wall clock is the axis on which a scheduler observes a job.

\textbf{Overhead.} The model's run takes $0.99\times$ the time of a plain RSS run at the median, and $0.96\times$ to
$1.04\times$ on 20 of the 23 programs. The plain RSS run includes the 40\,ms that the external sampler takes to start,
which explains the ratios of $0.60\times$ for atax and $0.84\times$ for NPB-MG, whose runs last 0.10\,s and 0.31\,s.
darknet runs $1.21\times$ slower under the model in all 3 runs. A thread that sleeps for 0.99\,s between samples slows
darknet by the same $1.22\times$, so the cost comes from the presence of a second thread in a single-threaded program,
and not from the sampling. We did not identify which code path the second thread slows down. The mode that reads
\texttt{smaps} takes $0.98\times$ the time of a plain RSS run at the median ($0.60\times$ to $1.18\times$).
\memprint{}'s W5 run with its counter run takes $12.0\times$ a native run at the median, $8.9\times$ on the other
programs and $34.4\times$ on the allocator-heavy programs, with $93.0\times$ on \texttt{cc1}.

\subsection{Where the remaining RSS goes}
\label{sec:gap}

\begin{table}[t]
\centering
\small
\caption{Parts of the model's run, each summed over the run as a percentage of the truth's sum (the truth is 100).
Data is estimate 1. Free is glibc's free bytes. File RSS and file PSS are the RSS and PSS of file-backed pages (RSS
minus anonymous RSS, and PSS minus anonymous RSS). Median of 3 runs.}
\label{tab:components}
\begin{tabular}{lrrrr}
\toprule
Program & Data (\%) & Free (\%) & File RSS (\%) & File PSS (\%) \\
\midrule
miniVite & 92 & 38 & 87 & 56\\
python & 101 & 15 & 21 & 9\\
lua & 101 & 23 & 7 & 1\\
sqlite & 108 & 7 & 9 & 4\\
perl & 95 & 15 & 10 & 4\\
quicksilver & 96 & 2 & 23 & 15\\
cc1 & 93 & 2 & 23 & 20\\
\midrule
others (median of 16) & 99 & 1 & 5 & 0\\

\bottomrule
\end{tabular}
\end{table}

\tabref{tab:components} shows the size of each part of the model's run against the truth. On the allocator-heavy
programs, glibc holds free memory equal to 6\% to 39\% of the truth on 5 of the 7 programs, and file-backed pages make
up 6\% to 88\% of the truth. On the other programs, both parts are small, at a median of 1\% and 6\%.

miniVite has the largest share of both parts. glibc's free bytes amount to 39\% of the truth, and file-backed pages
amount to 88\%. A sample of \texttt{/proc/pid/smaps} in the middle of a run shows where the file-backed pages are. The
MPI library holds 2348\,KiB, PMIx 2096\,KiB, the C library 1496\,KiB and the C++ runtime 1356\,KiB, against a truth of
9.7\,MiB on average over the run. The PSS of the C library is 77\,KiB and that of the C++ runtime 199\,KiB, because
other processes map the same pages, but the PSS of the MPI and PMIx libraries equals their RSS, because only miniVite
maps them. Subtracting glibc's free bytes alone therefore reduces miniVite's error only from $+105\%$ to $+66\%$
(\secref{sec:options}), and leaving out the libraries' pages reduces it to $-7\%$.

\section{Transfer to a larger input}
\label{sec:transfer}

\begin{table}[t]
\centering
\small
\caption{Whole-run error (\% of the truth), signed and absolute, of each estimate predicted for a program's largest
input from its smaller inputs. The peak is a linear regression of the smaller inputs' peaks on the input size, never
below the largest of them, and the shape is the largest smaller input's curve. -- marks a program whose smaller inputs
show nothing above the runtime's pages, for which the scheduler falls back to the user's request.}
\label{tab:transfer}
\setlength{\tabcolsep}{5pt}
\begin{tabular}{lrrrrrrrr}
\toprule
 & \multicolumn{2}{c}{RSS} & \multicolumn{2}{c}{Data (1)} & \multicolumn{2}{c}{Shared (2)} & \multicolumn{2}{c}{\memprint{}} \\
\cmidrule(lr){2-3}\cmidrule(lr){4-5}\cmidrule(lr){6-7}\cmidrule(lr){8-9}
Program & sign. & abs. & sign. & abs. & sign. & abs. & sign. & abs. \\
\midrule
\tabinput{tables/transfer.tex}
\bottomrule
\end{tabular}
\end{table}

A scheduler rarely has a profile of the job's own input, so a footprint model is more useful if a curve measured on
small inputs predicts the curve of a larger one. We apply one transfer rule, fixed before looking at the results, to
every estimate. The peak at the largest input is a linear regression of the smaller inputs' peaks on the input
size~\cite{KSplus2024}, never below the largest peak seen, and the shape is the curve of the largest smaller input,
scaled to that peak. The largest input is held out, and the model's smaller inputs run once each.

The data estimate transfers as accurately as \memprint{} (\tabref{tab:transfer}). On the allocator-heavy programs, its
median whole-run error is $-4\%$ (10\% absolute), against $+20\%$ (22\%) for RSS and $-6\%$ (10\%) for \memprint{}. On
the other programs, it is $-2\%$ (6\%), against $+2\%$ (10\%) for RSS and 0\% (5\%) for \memprint{}.

Three cases show where transfer fails. On \texttt{cc1}, every estimate is 226\% to 351\% too high, because the three
inputs are different source files whose size predicts the compiler's memory poorly. On NPB-MG, the data estimate is
67\% too low and RSS 96\% too low, because the only usable smaller input runs for a few milliseconds and the 5\,ms
sampling misses its peak, while \memprint{} counts every reference and is 15\% too high. On 2mm, gemm and jacobi-2d,
RSS is 25\% to 35\% too low for the same reason, because the smallest input's run fits in one sample of the external
sampler, which sees only the runtime. The model's thread starts before \texttt{main} and samples these runs, and its
error on them is $-10\%$ to $-4\%$. The transfer of a model trained on short runs therefore depends on sampling the
start of a run, and sampling faster during the first milliseconds of a run would address it.

\section{Scheduling with the model's curves}
\label{sec:sched}

\begin{table}[t]
\centering
\small
\caption{Scheduling simulation under the pages-touched truth on the native clock: mean response time (minutes) at
three offered loads, and throughput (jobs per hour), OOM kills and core-hours lost to them at a load of 0.8. Means
over 3 seeds of 3000 jobs on 32 nodes. Every estimate's margin is the 90th percentile of the margins the other
programs need (leave one program out), and the user's request takes 10\%. The second block profiles the job's own
input, and the third transfers from the program's smaller inputs.}
\label{tab:sched}
\setlength{\tabcolsep}{5pt}
\begin{tabular}{lrrrrrr}
\toprule
 & \multicolumn{3}{c}{Mean response (min)} & \multicolumn{3}{c}{At load 0.8} \\
\cmidrule(lr){2-4}\cmidrule(lr){5-7}
Oracle & 0.6 & 0.8 & 1.0 & jobs/h & kills & lost core-h \\
\midrule
Exact curve (truth) & 55 & 94 & 261 & 62.3 & 0 & 0\\
Exact peak & 172 & 479 & 698 & 47.4 & 0 & 0\\
User request & 1867 & 2244 & 2452 & 23.6 & 0 & 0\\
KS+ & 491 & 976 & 1233 & 33.7 & 9236 & 5012\\
\midrule
RSS peak & 532 & 875 & 1126 & 35.4 & 412 & 886\\
RSS curve & 185 & 515 & 776 & 41.3 & 407 & 1732\\
\memprint{} W5 & 1247 & 1607 & 1797 & 27.2 & 386 & 1363\\
\memprint{} W5, native clock & 197 & 519 & 707 & 42.6 & 384 & 593\\
Data $\hat f_1$ & 159 & 433 & 644 & 44.9 & 319 & 1888\\
Charged $\hat f_2$ & 192 & 508 & 691 & 41.7 & 332 & 1916\\
\midrule
RSS peak & 1913 & 2269 & 2476 & 22.3 & 418 & 7\\
RSS curve & 2369 & 2632 & 2869 & 20.3 & 698 & 553\\
\memprint{} W5 & 1890 & 2273 & 2475 & 22.3 & 387 & 1363\\
\memprint{} W5, native clock & 911 & 1283 & 1477 & 30.5 & 275 & 699\\
Data $\hat f_1$ & 2523 & 2806 & 3142 & 20.4 & 926 & 570\\
Charged $\hat f_2$ & 2017 & 2479 & 2665 & 21.0 & 939 & 1554\\

\bottomrule
\end{tabular}
\end{table}

We place the model's curves in DMSched-Sim, the trace-driven scheduling simulator of the corpus. The simulator replays
a synthetic job stream with MetaCentrum's shape (3000 jobs, 1 to 16 cores, run times of 5 to 720 minutes) on 32 nodes
with 32 cores and 128\,GB each and a 1\,TB memory pool. Each job takes the footprint curve of a random corpus program
with a smaller input, scaled to the job's peak, and the scheduler reserves memory by the oracle's estimate of that
curve. The reservation divides the run into 20 windows, holds the estimate's largest value in each window, widens each
window to its neighbours, and adds a margin times the peak. A job whose memory exceeds its reservation is killed, loses
its work, and is requeued with a reservation 1.5 times larger. Each oracle's margin is calibrated by leaving one
program out, so a job of program $p$ gets the 90th percentile of the margins the other programs need to never fall
below their truth.

On the pages truth (\tabref{tab:sched}), the data estimate gives the lowest mean response time of all estimates that
profile the job's own input, 433 minutes at a load of 0.8, against 515 for the RSS curve, 519 for \memprint{} on the
native clock and 479 for the exact peak. It also has the highest throughput of these estimates (44.9 jobs per hour)
and the fewest kills (319). Its calibrated margin is 7\%, against 12\% for \memprint{}, because no program needs a
margin above 9\% with it. \memprint{} on Pin's own clock gives 1607 minutes, because its curve and the job's footprint
disagree on when memory is used.

Every transferred estimate except \memprint{} on Pin's clock gives a response time at least twice that of the same
estimate on the job's own input. \memprint{} on the native clock transfers best (1283 minutes), and the data estimate
transfers worst (2806 minutes), although its median transfer error equals \memprint{}'s (\secref{sec:transfer}). The
calibrated margin follows the worst programs, and the data estimate needs a margin of 1.00 on gemm and 2.67 on NPB-MG,
the programs whose smaller inputs are too short for 5\,ms sampling, so its 90th percentile is 0.53, against 0.23 for
\memprint{}.

Under the RSS truth, which the kernel charges and OOM-kills on, the RSS curve of an earlier run is the best estimate
(393 minutes at a load of 0.8), because it measures the truth itself. The data estimate gives 748 minutes, close to
\memprint{} on the native clock (699 minutes), because it leaves out the retained memory and libraries that the RSS
truth counts, and its calibrated margin must cover them. The estimate a scheduler should use therefore depends on what
the node enforces. A limit on RSS calls for the RSS curve, and a limit on the memory the job uses, as in a cluster
whose allocators return freed memory, calls for the data estimate.

\section{Other allocators and transparent huge pages}
\label{sec:alloc}

\begin{table}[t]
\centering
\small
\caption{Signed whole-run error (\% of the truth under the same allocator) under jemalloc and TCMalloc, with
transparent huge pages enabled (the machine's setting) and disabled for the process. \emph{Anon} is the data estimate
without the allocator's free bytes. Median of 3 runs. Overhead is the run time over a plain RSS run under the same
allocator.}
\label{tab:alloc}
\setlength{\tabcolsep}{5pt}
\begin{tabular}{lrrrrrrr}
\toprule
 & \multicolumn{4}{c}{THP enabled} & \multicolumn{3}{c}{THP disabled} \\
\cmidrule(lr){2-5}\cmidrule(lr){6-8}
Program & RSS & Anon & Data (1) & Shared (2) & RSS & Data (1) & Shared (2) \\
\midrule
\tabinput{tables/alloc.tex}
\bottomrule
\end{tabular}
\end{table}

\begin{table}[t]
\centering
\small
\caption{The corpus under glibc with transparent huge pages enabled and disabled for the process: median signed and
absolute whole-run error (\% of the truth). The data estimate here comes from the mode that reads \texttt{smaps}.}
\label{tab:thp}
\begin{tabular}{lrrrrrrrr}
\toprule
 & \multicolumn{4}{c}{RSS} & \multicolumn{4}{c}{Data (1)} \\
\cmidrule(lr){2-5}\cmidrule(lr){6-9}
 & \multicolumn{2}{c}{THP on} & \multicolumn{2}{c}{THP off} & \multicolumn{2}{c}{THP on} & \multicolumn{2}{c}{THP off} \\
Programs & sign. & abs. & sign. & abs. & sign. & abs. & sign. & abs. \\
\midrule
Allocator-heavy (7) & +16 & 17 & +15 & 15 & $-$3 & 9 & $-$7 & 9\\
Others (16) & +1 & 2 & 0 & 3 & $-$1 & 3 & $-$1 & 4\\

\bottomrule
\end{tabular}
\end{table}

We ran the 7 allocator-heavy programs under jemalloc 5 and TCMalloc (gperftools 2.18, \texttt{tcmalloc\_minimal}),
loaded with \texttt{LD\_PRELOAD}, and measured the truth with Pin under the same allocator (\tabref{tab:alloc}).

\textbf{TCMalloc.} RSS overestimates the footprint by a median of 46\%, more than under glibc, because TCMalloc keeps
freed spans in its page heap and its caches. The data estimate has a median error of $-4\%$, and it is low on SQLite
($-22\%$) and miniVite ($-16\%$). Without the allocator's free bytes, the estimate is $+21\%$, so the statistics remove
most of the gap.

\textbf{jemalloc.} With transparent huge pages enabled, every estimate overestimates jemalloc's footprint, RSS by a
median of 56\% and the data estimate by 35\%, with $+149\%$ on miniVite. jemalloc places its extents at 2\,MiB aligned
addresses, so the kernel backs them with huge pages, and a 2\,MiB page counts in full when the program touches 4\,KiB
of it. In a run of Lua under jemalloc, 45\,MiB of the 65\,MiB of anonymous memory is in huge pages. We disabled
transparent huge pages for the process with \texttt{prctl(PR\_SET\_THP\_DISABLE)}, which the library calls in its
constructor on request. The data estimate's median error under jemalloc then drops to $-1\%$, and RSS's to $+26\%$.
Under TCMalloc and under glibc, disabling transparent huge pages changes the median errors by at most 4 points
(\tabref{tab:thp}).

The memory that huge pages add is real. A node with transparent huge pages enabled holds the whole 2\,MiB page in
memory, and the pages-touched truth undercounts it, as \memprint{} would. A scheduler on such a node needs the estimate
with huge pages, which the model measures and Pin does not, and a scheduler that wants the 4\,KiB footprint can disable
huge pages for the profiling run. We also checked two other kernel counters that we expected to matter.
\texttt{LazyFree}, the pages released with \texttt{MADV\_FREE} that stay resident until the kernel reclaims them, was 0
in the jemalloc runs, and the \texttt{Referenced} bytes of llama2.c's model file equal its RSS.

\section{Language model engines}
\label{sec:llm}

\begin{table}[t]
\centering
\small
\caption{Signed whole-run error (\% of the pages-touched truth) on the language model engines, median of 3 runs. Data
(1) without files is $\max(A - F, 0) + B$. The model's overhead is the run time of the mode that reads \texttt{smaps}
over a plain RSS run, and \memprint{}'s is its W5 and counter runs over a native run ($\times$).}
\label{tab:llm}
\setlength{\tabcolsep}{4.5pt}
\begin{tabular}{lrrrrrrrr}
\toprule
 & & \multicolumn{2}{c}{Data (1)} & & & Data (1) & \multicolumn{2}{c}{Overhead ($\times$)} \\
\cmidrule(lr){3-4}\cmidrule(lr){8-9}
Engine and input & RSS & no files & with files & Shared (2) & \memprint{} & THP off & model & \memprint{} \\
\midrule
llama2c 110M & +12 & $-$86 & +12 & +11 & 0 & 0 & 1.00 & 3.9\\
llama2c 15M & 0 & $-$92 & 0 & $-$3 & $-$2 & 0 & 0.98 & 6.0\\
llama2c 42M & +15 & $-$81 & +15 & +14 & 0 & $-$1 & 1.00 & 4.4\\
llamacpp 135m-n128 & +1 & $-$39 & $-$1 & 0 & $-$3 & $-$7 & 0.99 & 38.9\\
llamacpp 135m-n32 & $-$26 & $-$56 & $-$28 & $-$27 & $-$13 & $-$30 & 0.94 & 67.4\\

\bottomrule
\end{tabular}
\end{table}

llama2.c maps its model file and reads the weights from the mapping, and llama.cpp maps its GGUF file and copies part
of it into buffers it allocates. Without the data-file term, the data estimate is 81\% to 92\% too low on llama2.c and
39\% to 56\% too low on llama.cpp (\tabref{tab:llm}), because the weights are file-backed pages. With the term, the
data estimate is within $-1\%$ to $+15\%$ of the truth on llama2.c and on llama.cpp generating 128 tokens. On
llama2.c, the remaining overestimate comes from transparent huge pages. The engine allocates its key and value caches
for the whole context of 1024 positions with \texttt{calloc} and fills only the positions it generates, and 45\,MiB of
its 47.5\,MiB of anonymous memory is in huge pages. With huge pages disabled, the data estimate is within 1\% of the
truth on all three model sizes.

llama.cpp generating 32 tokens is the exception. Every estimate of the model is 26\% to 28\% too low, against 13\% for
\memprint{}. The run lasts less than a second, and most of it loads the model, so the estimates and the truth disagree
on when the model's pages become resident. The model's run costs $0.94\times$ to $1.00\times$ a plain RSS run, and
\memprint{}'s W5 costs $3.9\times$ to $6.0\times$ on llama2.c and $38.9\times$ to $67.4\times$ on llama.cpp.

\section{Other ways to correct RSS}
\label{sec:options}

\begin{table}[t]
\centering
\small
\caption{Signed whole-run error (\% of the truth) of each alternative used alone on the allocator-heavy programs,
median of 3 runs, and its overhead over a plain RSS run ($\times$). Runtime statistics exist only for CPython and Lua.
These runs keep the file-backed pages in every estimate.}
\label{tab:options}
\setlength{\tabcolsep}{5pt}
\begin{tabular}{lrrrrrr}
\toprule
Program & RSS & RSS$-$free & trim 5\,ms & trim thr.\ 0 & runtime stats & referenced \\
\midrule
miniVite & +103 & +64 & +56 & +53 & -- & $-$70\\
python & +30 & +14 & +15 & +16 & +9 & +15\\
lua & +23 & +2 & +5 & +19 & $-$18 & +4\\
sqlite & +17 & +12 & +15 & +18 & -- & $-$18\\
perl & +16 & +1 & +10 & +12 & -- & $-$12\\
quicksilver & +16 & +16 & +13 & +13 & -- & $-$56\\
cc1 & +14 & +12 & +12 & +13 & -- & $-$69\\
\midrule
median signed & +17 & +12 & +13 & +16 & $-$4 & $-$18\\
median absolute & 20 & 14 & 15 & 17 & 38 & 32\\
overhead ($\times$) & 1.00 & 0.99 & 1.03 & 1.02 & 0.96 & 1.02\\

\bottomrule
\end{tabular}
\end{table}

Before we settled on the model, we tested four alternatives on the allocator-heavy programs under glibc, each on its
own (\tabref{tab:options}). The first, RSS minus glibc's free bytes, is the first half of the model.

\textbf{Releasing freed memory during the run.} Instead of subtracting free bytes, the program can return them to the
kernel, so that RSS drops when the program frees memory. We tested two ways. The first calls \texttt{malloc\_trim(0)}
every 5\,ms from a thread in the process, which releases the free pages inside every arena with \texttt{madvise}. The
second sets glibc's trim threshold to 0, which releases free memory at the top of the heap after every \texttt{free}.
Both reduce the median error less than subtracting free bytes, $+11\%$ and $+14\%$ against $+9\%$. A trim threshold of
0 releases only the top of the heap, which explains why it changes Lua's error by only 3 points. Both alternatives
also change the program's allocation behaviour, and the trim thread costs up to $1.10\times$ on Perl.

\textbf{Language runtime statistics.} For CPython, we sampled the free bytes of its object allocator (the arena bytes
minus the bytes in allocated blocks, from \texttt{sys.\_debugmallocstats}) from a thread started in
\texttt{sitecustomize}. For Lua, we sampled the heap size the garbage collector reports
(\texttt{collectgarbage("count")}) from a debug hook every $10^6$ virtual machine instructions. CPython's whole-run
signed error drops to $+5\%$, but its absolute error rises, which means that the estimate is too low in some parts of
the run and too high in others. We did not determine why. Lua's heap count leaves out the interpreter's own memory and
gives $-19\%$. The statistics exist only for these two programs, and GCC reports its heap only at the end of a
compilation.

\textbf{Referenced pages per window.} The kernel can clear the referenced bits of a process's pages
(\texttt{/proc/pid/clear\_refs}), after which \texttt{smaps\_rollup} counts the pages the process has touched since. We
cleared the bits 20 times per run and took the pages referenced in each window as the footprint of that window, in the
manner of a working set~\cite{denning1968}. The estimate falls below the truth by up to 70\% (miniVite), 68\%
(\texttt{cc1}) and 55\% (Quicksilver), because a page that the program holds but does not touch within a window is not
counted. A scheduler that reserved this estimate would run these programs out of memory.

\section{Returning free memory at run time}
\label{sec:trim}

The model's free bytes are memory that the program has freed and that the allocator keeps resident. A runtime that
reads the same statistics can return this memory to the kernel while the program runs, which lowers the RSS and the
memory that a system must provision for the job. We added the model to libmpr, the runtime of our allocation-resizing
study, which serves each allocation of at least 64\,KiB from a per-site reservation that \memprint{} profiles and
commits only the predicted part of the block. libmpr leaves smaller blocks to glibc, so these blocks hold the free
bytes that the model subtracts.

\textbf{Mechanism.} With a share $s$ set (\texttt{MPR\_TRIM}), libmpr compares glibc's resident free bytes with
$\max(s \cdot A, 1\,\mathrm{MB})$, where $A$ is the anonymous RSS that the kernel keeps as a counter in
\texttt{/proc/self/status}, and it calls \texttt{malloc\_trim(0)} when the free bytes exceed this threshold.
\texttt{malloc\_trim} releases the whole pages inside every free chunk of every arena with \texttt{madvise}. glibc
counts the released bytes as free after the trim, because they are free but no longer resident, so libmpr keeps the
free bytes at the last trim apart as released, and it lowers this count when the free bytes fall below it, which
means that the program has reused released memory. The resident free bytes are then the free bytes minus the
released bytes.

libmpr makes this check from \texttt{free}, which it already intercepts, and not from a thread. A second thread
turns off the paths that glibc takes in a single-threaded process, in which \texttt{stdio} and \texttt{malloc} skip
their locks. With a monitoring thread that never trims, darknet, which parses its files one character at a time,
runs $1.23\times$ slower, and with the check in \texttt{free} it runs at $0.96\times$. The resident free bytes grow
only by the bytes that the program frees, so \texttt{free} counts these bytes and makes the next check only after the
program has freed the distance between the resident free bytes and the threshold, and at least 1\,MB. A program that
frees memory in a short burst is therefore trimmed at the end of the burst and not at the next tick of a timer. A
thread is still available for the monitor that writes the model's estimate every 5\,ms, and it trims as well.

\textbf{Methodology.} We ran the 29 workloads of the resizing study: the 23 programs at their largest inputs, with
miniVite built with the resizing compiler pass, and six configurations of llama2.c and llama.cpp. Each workload ran
under six builds, glibc, libmpr with the reservations of its own input and libmpr with reservations profiled on
other inputs, each without and with trimming at $s = 0.1$, three times each, with RSS sampled every 20\,ms. Four
workloads ran at a time on the 32-core machine, which adds about 5\% of noise to each run time. Every run produced
the output of the glibc run.

\begin{table}[t]
\centering
\small
\caption{Memory summed over the workloads under each build, mean of 3 runs: the peak RSS and the time-averaged RSS
(MB), RSS integrated over the run (GB\,s), and the time-averaged reservation of dynamic provisioning, which holds the
request for the first sixth of the run and then each sixth's largest RSS (MB). Trimming at $s = 0.1$, from
\texttt{free}.}
\label{tab:trim-totals}
\setlength{\tabcolsep}{4pt}
\begin{tabular}{lrrrrrrrr}
\toprule
 & \multicolumn{4}{c}{Allocator-heavy (7)} & \multicolumn{4}{c}{All (29)} \\
\cmidrule(lr){2-5} \cmidrule(lr){6-9}
Build & Peak & Mean & RSS$\times$time & Dynamic & Peak & Mean & RSS$\times$time & Dynamic \\
 & (MB) & (MB) & (GB\,s) & (MB) & (MB) & (MB) & (GB\,s) & (MB) \\
\midrule
glibc & 544 & 360 & 0.9 & 478 & 4{,}590 & 3{,}600 & 17.1 & 4{,}390 \\
glibc + trim & 534 & 344 & 0.9 & 466 & 4{,}597 & 3{,}524 & 17.3 & 4{,}354 \\
libmpr & 539 & 361 & 1.0 & 474 & 3{,}764 & 2{,}979 & 15.4 & 3{,}549 \\
libmpr + trim & 530 & 345 & 0.9 & 466 & 3{,}749 & 2{,}959 & 15.6 & 3{,}532 \\
libmpr, held-out profile & 543 & 363 & 1.0 & 478 & 3{,}776 & 2{,}988 & 15.9 & 3{,}601 \\
libmpr, held-out + trim & 533 & 344 & 0.9 & 465 & 3{,}766 & 2{,}959 & 15.8 & 3{,}588 \\
\bottomrule
\end{tabular}

\end{table}

\begin{table}[t]
\centering
\small
\caption{The workloads on which trimming happened (more than one trim in a run): time-averaged RSS (MB), run time
over the glibc run's, and the trims per run, mean of 3 runs. On the other 17 workloads, libmpr made at most one
trim per run.}
\label{tab:trim}
\setlength{\tabcolsep}{4pt}
\begin{tabular}{lrrrrrrrr}
\toprule
 & \multicolumn{4}{c}{Time-averaged RSS (MB)} & \multicolumn{2}{c}{Run time ($\times$ glibc)} & \multicolumn{2}{c}{Trims} \\
\cmidrule(lr){2-5} \cmidrule(lr){6-7} \cmidrule(lr){8-9}
Workload & glibc & +trim & libmpr & +trim & +trim & libmpr+trim & glibc & libmpr \\
\midrule
miniVite (pass) & 23 & 18 & 20 & 19 & 1.03 & 1.05 & 10 & 7 \\
python & 52 & 48 & 53 & 48 & 1.01 & 1.01 & 9 & 9 \\
lua & 45 & 39 & 47 & 39 & 0.98 & 1.07 & 39 & 38 \\
sqlite & 42 & 41 & 43 & 41 & 1.06 & 1.10 & 26 & 24 \\
perl & 52 & 51 & 52 & 52 & 1.09 & 1.07 & 7 & 6 \\
quicksilver & 53 & 52 & 52 & 52 & 1.01 & 0.99 & 20 & 3 \\
\midrule
gap-sssp & 66 & 67 & 67 & 69 & 1.00 & 1.05 & 4 & 1 \\
llamacpp-135m-n128 & 294 & 291 & 141 & 140 & 1.05 & 1.05 & 20 & 18 \\
llamacpp-135m-n32 & 203 & 178 & 94 & 89 & 1.10 & 1.00 & 20 & 18 \\
llamacpp-360m-n64 & 624 & 583 & 335 & 335 & 0.97 & 0.99 & 20 & 18 \\
lulesh & 96 & 95 & 95 & 93 & 1.01 & 1.06 & 116 & 96 \\
minife & 260 & 259 & 260 & 259 & 1.03 & 1.00 & 12 & 3 \\
\bottomrule
\end{tabular}

\end{table}

\textbf{Memory.} Trimming lowers the time-averaged RSS of the allocator-heavy programs from 360 to 344\,MB in total
(4.4\%), and that of all 29 workloads from 3{,}600 to 3{,}524\,MB under glibc (2.1\%) and from 2{,}979 to
2{,}959\,MB under libmpr (0.7\%) (\tabref{tab:trim-totals}). The largest reductions are on miniVite (23 to 18\,MB),
Lua (45 to 39\,MB), CPython (52 to 48\,MB) and llama.cpp under glibc (203 to 178\,MB at 32 tokens)
(\tabref{tab:trim}). Under libmpr, llama.cpp gains less from trimming, because libmpr already serves its large
buffers and glibc holds less free memory. Trimming does not lower the summed peak RSS, 4{,}590 against 4{,}597\,MB,
because the peak of most programs comes before they free memory.

Trimming returns less memory than the model subtracts on some programs, because \texttt{malloc\_trim} can release
only whole pages inside free chunks. Perl holds 6.9\,MB of free bytes on average over its run, but trimming lowers
its time-averaged RSS by 0.3\,MB, which means that Perl's free bytes lie in chunks smaller than a page or in pages
that also hold live data. On Lua, trimming returns 6.1 of 8.3\,MB of free bytes, and on miniVite 4.8\,MB, as much as
the 4.4\,MB of free bytes of its build without the pass. On \texttt{cc1},
GCC's garbage-collected heap keeps its own free pages, and libmpr made no trim.

\textbf{Overhead.} The median run time with trimming is $1.01\times$ the glibc run's, and $1.02\times$ under libmpr
against $1.00\times$ for libmpr alone. The highest ratios are $1.10\times$, on SQLite under libmpr with 24 trims per
run and on llama.cpp at 32 tokens under glibc, but libmpr without trimming reaches $1.11\times$ on Perl, so these
differences are within the noise of our runs. The RSS integrated over the run does not fall with
trimming, 17.1 against 17.3\,GB\,s over all workloads, because the run times vary more than the RSS falls.

\begin{table}[t]
\centering
\small
\caption{Trimming under glibc at three shares $s$: the trims per run summed over the allocator-heavy programs, the
time-averaged RSS summed over the allocator-heavy programs and over all 29 workloads (MB), and the run time over the
glibc run's in the same set of runs ($\times$), median and largest over the workloads, mean of 3 runs. The untrimmed
row and the row at $s = 0.1$ come from the same runs.}
\label{tab:trim-share}
\setlength{\tabcolsep}{4pt}
\begin{tabular}{lrrrrr}
\toprule
 & Trims & \multicolumn{2}{c}{Time-averaged RSS (MB)} & \multicolumn{2}{c}{Run time ($\times$ glibc)} \\
\cmidrule(lr){3-4} \cmidrule(lr){5-6}
Share $s$ & per run & Alloc.-heavy & All & Median & Largest \\
\midrule
none & 0 & 360 & 3{,}600 & 1.00 & 1.00 \\
0.02 & 321 & 342 & 3{,}550 & 1.01 & 1.09 \\
0.1 & 111 & 344 & 3{,}524 & 1.01 & 1.10 \\
0.3 & 38 & 350 & 3{,}555 & 1.00 & 1.08 \\
\bottomrule
\end{tabular}

\end{table}

\textbf{Share.} We repeated the glibc runs with trimming at $s = 0.02$ and $s = 0.3$ (\tabref{tab:trim-share}). At
$s = 0.02$, libmpr trims 2.9 times as often as at $s = 0.1$, 321 against 111 trims per run on the allocator-heavy
programs, but the time-averaged RSS of these programs stays at 342 against 344\,MB. At $s = 0.3$, it trims 38 times
and the RSS falls by half as much, to 350\,MB. The memory that trimming returns is therefore limited by the whole free
pages that glibc holds and not by how often libmpr checks, and $s = 0.1$ returns nearly all of them. The total over
all 29 workloads varies by 30\,MB between the sets of runs without a trend in $s$, because llama.cpp's RSS varies
from run to run. The largest run time is $1.09\times$ on Perl at $s = 0.02$, with 40 trims per run.

\textbf{Returning every free byte.} To bound what any allocator could gain by returning its free memory, we also
subtracted all of glibc's free bytes from the RSS of the untrimmed runs, which is the RSS that an allocator would
hold if it returned every free byte to the kernel at once, as a compacting allocator such as Mesh~\cite{Powers2019Mesh}
aims to do. The free bytes come from libmpr's monitor in separate runs without trimming, placed on the same steps of
progress. glibc's \texttt{mallinfo} counts the free bytes of the main arena only, so the bound misses the free memory
of the arenas of other threads, which \texttt{malloc\_trim} releases. On llama.cpp, the bound is therefore above the
trimmed RSS, and we take the smaller of the two at each step. On the simulator's curves, from which the runtime's
pages are removed, the time-averaged RSS of the allocator-heavy programs falls from 352\,MB under glibc to 335\,MB
with trimming and to 318\,MB at the bound, and their summed peak RSS falls from 537 to 526 and 514\,MB. Trimming
therefore returns about half of the free memory, and the other half lies in free chunks smaller than a page or in
pages that also hold live data, most of it Perl's (51 against 44\,MB at the bound).

\begin{table}[t]
\centering
\small
\caption{Cluster simulation of the resizing study (DMSched-Sim, 3{,}000 jobs, the RSS as the job's truth): throughput
and mean response time of each policy with glibc's RSS, with trimming at $s = 0.1$ and at the bound of returning
every free byte, mean of 3 seeds, and the number of the 6 seed and load pairs on which the throughput is higher than
with glibc's RSS. The standard deviation of the response over seeds is about 200\,min for the dynamic policies.}
\label{tab:trim-grid}
\setlength{\tabcolsep}{4pt}
\begin{tabular}{lrrrrrrrr}
\toprule
 & \multicolumn{4}{c}{Load 0.8} & \multicolumn{4}{c}{Load 1.2} \\
\cmidrule(lr){2-5} \cmidrule(lr){6-9}
 & \multicolumn{2}{c}{Jobs/h} & \multicolumn{2}{c}{Response (min)} & \multicolumn{2}{c}{Jobs/h} & \multicolumn{2}{c}{Response (min)} \\
Policy & no trim & trim & no trim & trim & no trim & trim & no trim & trim \\
\midrule
Request held & 21.7 & 21.7 & 1{,}844 & 1{,}842 & 21.7 & 21.7 & 2{,}421 & 2{,}438 \\
Dynamic provisioning & 21.5 & 20.9 & 1{,}149 & 1{,}109 & 21.0 & 21.4 & 1{,}768 & 1{,}629 \\
Dynamic, look-back & 21.6 & 22.4 & 1{,}129 & 1{,}041 & 21.8 & 22.2 & 1{,}605 & 1{,}557 \\
libmpr & 25.0 & 25.3 & 1{,}246 & 1{,}194 & 25.1 & 25.4 & 1{,}817 & 1{,}797 \\
libmpr + dynamic & 23.1 & 24.4 & 775 & 705 & 23.2 & 24.2 & 1{,}340 & 1{,}350 \\
libmpr + dynamic, look-back & 24.1 & 24.4 & 741 & 708 & 25.2 & 24.5 & 1{,}242 & 1{,}279 \\
\bottomrule
\end{tabular}

\end{table}

\textbf{Cluster.} We repeated the cluster simulation of the resizing study with the trimmed builds and with the
bound, in which each job is one of the 29 workloads and its demand over time is its RSS (\tabref{tab:trim-grid}). A
job that holds its request for the whole run gains nothing from trimming, because the request does not change. Under
dynamic provisioning with a look-back policy, which resizes the reservation to the RSS of the previous interval,
trimming lowers the mean response time at load 0.8 from 1{,}129 to 1{,}041\,min, and the bound lowers it to
989\,min. With libmpr, trimming raises the throughput on every seed and load, but only from 25.0 to 25.3 jobs per hour
at load 0.8, and with libmpr and look-back dynamic provisioning, trimming raises it on 3 of 6, and the mean moves from 24.1
to 24.4 at load 0.8 and from 25.2 to 24.5 at load 1.2. The bound raises the throughput of every dynamic policy on every seed
and load except one, and with libmpr and look-back dynamic provisioning from 24.1 to 26.4 jobs per hour and the
response from 741 to 559\,min at load 0.8. The simulation does not charge the run time of trimming, and the bound has
no run time, because no allocator we ran achieves it.

We conclude that trimming returns a few megabytes per allocator-heavy program at no measurable cost, but that on
this corpus its effect on the cluster is within the variation over seeds, because libmpr's reservations on the large
blocks account for most of the memory that the jobs request and do not use. An allocator that returned the free
memory in chunks smaller than a page would raise the throughput of a cluster that already resizes the large blocks and the reservations by 0.9
to 2.3 jobs per hour.

\section{Costs and effects of the approaches}
\label{sec:space}

\begin{table}[!htbp]
\centering
\footnotesize
\caption{The approaches to the gap between RSS and the footprint that we measured, in this paper and in the resizing
study, grouped by where they act. One-time cost: paid once per program and input or model, as a run time over the
native run's ($\times$). Cost per run: the run time of every production run over glibc's ($\times$), median and the
largest we observed. Effect: on the 7 allocator-heavy programs unless stated (whole-run error against the truth, or
memory), and in the cluster simulation at load 0.8 (jobs per hour and mean response time).}
\label{tab:space}
\setlength{\tabcolsep}{3pt}
\begin{tabular}{L{3.1cm}L{2.2cm}L{2.2cm}L{4.3cm}L{3.6cm}}
\toprule
Approach & One-time cost & Cost per run & Effect & Limit \\
\midrule
\multicolumn{5}{l}{\emph{Estimate the footprint for the scheduler (the job's memory does not change)}} \\
RSS of a native run & $1.00\times$ & none & $+16\%$ whole-run error & counts free memory and libraries \\
\memprint{} W5 (Pin) & $12.0\times$, $18.7\times$ to $93.0\times$ & none & $-5\%$; per-site attribution & Pin's cost
on allocator-heavy programs \\
Data estimate (\secref{sec:model}) & $0.99\times$, $1.21\times$ darknet & none & $-4\%$; $-4\%$ transferred;
433 against 515\,min for the RSS curve (\secref{sec:sched}) & glibc, jemalloc or TCMalloc only; THP; runs of a few ms \\
Runtime statistics & $0.96\times$ & none & $+5\%$ CPython, $-19\%$ Lua & two programs \\
Referenced pages & $1.02\times$ & none & $-24\%$, $-70\%$ miniVite & below the footprint: unsafe \\
\addlinespace
\multicolumn{5}{l}{\emph{Lower the memory the job holds (the runtime changes)}} \\
glibc trim threshold 0 & none & $1.02\times$ & RSS error $+17\%$ to $+14\%$ & top of the heap only \\
\texttt{malloc\_trim}, thread every 5\,ms & none & $1.03\times$, $1.23\times$ darknet & RSS error $+17\%$ to $+11\%$ &
second thread slows single-threaded programs \\
\texttt{malloc\_trim} from \texttt{free} (\secref{sec:trim}) & none & $1.01\times$, $1.10\times$ & time-averaged RSS
$-4.4\%$, peak unchanged; jobs/h $+0.3$ with libmpr & whole free pages only \\
Every free byte returned (bound) & none & not measured & time-averaged RSS $-9.7\%$, peak $-4.3\%$; jobs/h $+0.7$
with libmpr, $+2.3$ with libmpr and look-back & needs an allocator that compacts \\
libmpr reservations & Pin $217\times$ median, or libmpr's profile, one native run & $0.92\times$ to $1.03\times$ & 23
programs: peak commit $-10.7\%$, RSS unchanged; 29 workloads: jobs/h 21.7 to 25.0 & blocks of 64\,KiB or more \\
libmpr \texttt{MAP\_NORESERVE} stacks & none & $0.87\times$ to $1.13\times$ & commit $-8$ to $-34$\,MB per multithreaded program &
commit only \\
\addlinespace
\multicolumn{5}{l}{\emph{Reserve less than the request (the scheduler changes)}} \\
Dynamic provisioning, look-back & none & resizes and restarts, not charged here & response 1{,}844 to 1{,}129\,min & a
job restarts when its growth does not fit \\
libmpr and look-back & as libmpr & as both & response 741\,min; 559\,min at the bound & \\
Exact curve (upper bound) & --- & --- & 30.5 jobs/h, 630\,min; 34.3 and 340\,min with libmpr & needs the future \\
\bottomrule
\end{tabular}
\end{table}

The approaches that we measured fall into three groups, which differ in what they change and in when they cost
(\tabref{tab:space}). The first group estimates the footprint for a scheduler and leaves the job's memory unchanged.
Its cost is paid once per program and input, and a production run costs nothing. The second group lowers the memory
that the job holds, by returning free memory to the kernel or by committing only the predicted part of each large
block. Its cost is paid on every run, and libmpr's reservations also need a profile. The third group changes the
scheduler, which reserves less than the request and resizes the reservation while the job runs, at the cost of
resizes and of restarts when a growth does not fit.

\textbf{Estimating the footprint.} The data estimate gives \memprint{}'s accuracy at the cost of a native run, and
the other cheap estimates are either less accurate (RSS, the trim alternatives) or below the footprint (referenced
pages), which a scheduler cannot use without a margin. Pin remains the only route to a per-site profile, and
\memprint{} costs $12.0\times$ at the median.

\textbf{Lowering the memory.} Each approach in this group acts on a different part of the waste. libmpr's
reservations act on the large blocks that the program requests and does not touch, which is where most of the
request's waste lies on this corpus: they raise the throughput from 21.7 to 25.0 jobs per hour without changing RSS.
Trimming acts on the small blocks that the program has freed. From \texttt{free}, it costs $1.01\times$ at the median
and returns half of the free memory, and the other half needs an allocator that can return memory in chunks smaller
than a page. The cost of such an allocator is the open question of this group. Mesh reports its own costs, and we did
not run it.

\textbf{Reserving less.} Dynamic provisioning gains the most in response time, 1{,}844 to 1{,}129\,min, and its
costs are those of the cluster, resizes and restarts, and not of the job. The approaches combine. libmpr with
look-back dynamic provisioning gives 741\,min, and the bound on returning free memory lowers it to 559\,min, which is
below the 630\,min of the exact curve without libmpr. The exact curve with libmpr, 340\,min, is the limit that a
predicted curve would approach.

\section{Limitations}
\label{sec:limits}

A limitation of the free-byte correction is that it reads only the statistics of glibc, jemalloc and TCMalloc, which
leaves the memory of other allocators in the estimate. The arenas of CPython's object allocator and GCC's
garbage-collected heap stay in the estimate as RSS. On our corpus, the data estimate is nevertheless within 5\% of the
truth on CPython and \texttt{cc1}, and it is unclear whether this result holds for inputs on which these allocators
retain more memory.

A limitation of the data-file rule is that it classifies mappings by path, which counts any mapped file that is not a
shared object as data. A program that maps a large file it uses as code, or a library without \texttt{.so} in its
name, would have that file counted.

A limitation of the free-byte correction is also that some free bytes may not be resident, such as the untouched part
of the top of the heap, which makes the data estimate lower than the truth. The data estimate is below the truth on 6
of the 7 allocator-heavy programs under glibc, by 1\% to 9\% over the run, and on SQLite under TCMalloc by 22\%, so a
scheduler that uses it needs a margin.

A limitation of the sampling thread is that it adds a second thread to a single-threaded program, which slows darknet
by $1.21\times$ (\secref{sec:repeat}). Sampling from an external process avoids the thread, but the allocator's
statistics can be read only inside the process. libmpr's trimming avoids the thread by checking from \texttt{free}
(\secref{sec:trim}), but a monitor cannot, because it must sample while the program allocates without freeing. A second limitation of the thread is its 5\,ms interval, which misses
the curve of runs shorter than a few tens of milliseconds and limits the model's transfer from the smallest inputs
(\secref{sec:transfer} and \secref{sec:sched}).

The evaluation covers a single machine, short runs of 0.1 to 12\,s and single-process programs. The truth's clock
mapping passes through the native run's reference counts, which also place \memprint{}'s curve, so \memprint{}'s timing
accuracy is overstated against the truth. The truths under jemalloc and TCMalloc use the clock of the native run under
glibc.

\section{Related work}
\label{sec:related}

\textbf{Footprint measurement.} Valgrind's Massif and DHAT~\cite{Nethercote2007Valgrind,dhat} and \memprint{} measure
the heap or the touched memory of every allocation by instrumenting the program. Our model reads counters that the
kernel and the allocator already keep. Unlike \memprint{}, it does not attribute memory to allocation sites or show
which bytes of a block the program touches.

\textbf{Page access sampling.} DAMON~\cite{park2019damon} and Thermostat~\cite{agarwal2017thermostat} sample the
accessed bits of pages to find cold memory, and PEBS sampling~\cite{roca_nonell2020pebs} records the addresses of
sampled loads. These methods measure which pages a program uses within a period of time, which is the working
set~\cite{denning1968} and not the footprint. Our referenced-page alternative (\secref{sec:options}) shows that the
working set of a window falls far below the footprint on allocator-heavy programs.

\textbf{Memory prediction for scheduling.} KS+~\cite{KSplus2024} predicts a workflow task's memory over time from the
RSS of earlier runs, and Zacarias et al.~\cite{Zacarias2021Accurate,Zacarias2023Dynamic} predict the memory of HPC jobs
from earlier runs. Pond~\cite{Li2023Pond} predicts how much of a virtual machine's memory goes untouched. These systems
learn from RSS or from allocated memory, and our model gives them a training signal without the allocator's retained
memory and the job's shared libraries at the same cost as RSS.

\textbf{Allocators and huge pages.} Mesh~\cite{Powers2019Mesh} compacts the heap to return fragmented memory to the
kernel, and TCMalloc's Temeraire~\cite{Hunter2021Temeraire} releases free huge pages. Ingens~\cite{Kwon2016Ingens} and
HawkEye~\cite{Panwar2019HawkEye} limit the memory that huge pages waste on partly used ranges. These systems reduce the
gap between RSS and the footprint by changing the system, whereas our model measures the gap without changing the
program's allocations.

\section{Conclusion}
\label{sec:conclusion}

RSS overestimates the footprint of programs that allocate and free memory throughout their run, because the allocator
keeps freed memory resident and because RSS counts shared libraries. A model that subtracts the allocator's free bytes
from the anonymous RSS and adds mapped data files and the runtime's pages reduces the median whole-run error on these
programs from $+16\%$ to $-4\%$, which matches \memprint{}'s accuracy on the job's own input and on larger inputs, at a
median of $0.99\times$ the run time of a plain RSS run instead of $34.4\times$. In a scheduling simulation, its curve
gives the lowest response time of the estimates that profile the job's own input. Pin-based profiling remains needed
for attributing memory to code, for allocators the model cannot query, for the curves of runs too short to sample, and
for the 4\,KiB footprint of programs whose memory the kernel backs with huge pages. Used at run time, the same
statistics let a runtime return the allocator's free memory to the kernel, which lowers the RSS of allocator-heavy
programs by a few percent at no measurable cost, but which does not change the throughput of a cluster that already
resizes the large allocations. glibc returns only whole free pages, which is half of the free memory, and an allocator
that returned the rest would raise that cluster's throughput by 0.9 to 2.3 jobs per hour.

\textbf{Artifacts.} \texttt{paper\_rss/tools/mprshim.c} is the sampling library. \texttt{run\_model.sh},
\texttt{run\_transfer.sh}, \texttt{run\_options.sh}, \texttt{run\_alloc.sh} and \texttt{run\_llm.sh} run the
experiments, \texttt{grid.py} builds the simulator's corpus and runs the scheduling grid, and \texttt{compute.py}
produces every table. The runs are in \texttt{\textasciitilde/memprint-data/2026-10-26-rsscurve}.

\bibliographystyle{plain}
\bibliography{references}

\end{document}
